\section{Weighted Grassmann block integration}\label{sec:weighted-integration}

\subsection{Weighted Grassmann integration}\label{subsec:weighted-grassmann}

Let $\phi_i(t)$, $1\le i\le p$, be independent Grassmann families
on finite ordered index sets. If $\phi_i(I_i)$ is the ordered product
over $I_i$, write
\[
 F=\sum_{I_1,\ldots,I_p}c_F(I_1,\ldots,I_p)
                 \phi_1(I_1)\cdots\phi_p(I_p),\qquad
 \|F\|_{\mathsf w}=\sum_I|c_F(I)|\prod_i\mathsf w_i^{|I_i|}.
\]
The positive measures $\nu_i(t)$ on the index sets enter the kernels,
not this coefficient norm. Thus an antisymmetric kernel expansion
\[
 F=\sum_n\prod_i\frac1{n_i!}\int\dd t\,K_n(t)
                       \prod_i\prod_{j=1}^{n_i}\phi_i(t_j^i)
\]
has norm
\begin{equation}\label{eq:grassmann-kernel-norm}
 \|F\|_{\mathsf w}
     =\sum_n\prod_i\frac{\mathsf w_i^{n_i}}{n_i!}
                      \int\dd t\,|K_n(t)|.
\end{equation}
For arbitrary kernels the right side is an upper bound. Products of
ordered monomials vanish on repeated generators and otherwise multiply
their weights, so
\begin{equation}\label{eq:grassmann-product-bound}
 \|FG\|_{\mathsf w}\le\|F\|_{\mathsf w}\|G\|_{\mathsf w},
 \qquad \|e^F-1\|_{\mathsf w}\le e^{\|F\|_{\mathsf w}}-1
       \quad(F\text{ even}).
\end{equation}

For an additional finite nonempty mode set $\mathcal I$ with positive weights
$\nu(t)$, introduce $W(t),\bar W(t)$ and set
\[
 \langle\bar W,OW\rangle_\nu
    =\sum_{t,t'}\nu(t)\nu(t')\bar W(t)O(t,t')W(t'),
 \qquad \mathrm{Top}=\prod_{t\in\mathcal I}W(t)\bar W(t).
\]
The kernel represents $O$ with measure $\nu$; its matrix entry is
$O(t,t')\nu(t')$. Define $\int\mathrm DW$ as extraction of the
coefficient of $\mathrm{Top}$, multiplied by
$\prod_t\nu(t)^{-1}$, and define
\begin{align}
 \int\dd\mu_I(W)F
     &=\int\mathrm DW\,e^{-\langle\bar W,W\rangle_\nu}F,
                                            \label{eq:weighted-Berezin}\\
 \int\dd\mu_C(W)F
     &=(\det C)\int\mathrm DW\,
               e^{-\langle\bar W,C^{-1}W\rangle_\nu}F
                    \quad(C\text{ invertible}).\notag
\end{align}
These are algebra-valued functionals when other Grassmann generators
are present. The order of modes in $\mathrm{Top}$ is immaterial because
its factors are even.

\begin{proposition}\label{prop:weighted-Berezin-bounds}
Put $w_W=\max_t\nu(t)^{-1/2}$. Then
\begin{equation}\label{eq:weighted-integration-bound}
 \left\|\int\dd\mu_I(W)F\right\|_{\mathsf w}
                  \le\|F\|_{(\mathsf w,w_W)}.
\end{equation}
The integral factors for polynomials involving disjoint sets of modes.
For an invertible covariance,
\begin{align}
 \int\dd\mu_C(W)F(\bar W,W)
     &=\int\dd\mu_I(W)F(\bar W,CW),\label{eq:covariance-substitution}\\
 \int\dd\mu_C(W)W(t)\bar W(t')&=C(t,t'),\notag\\
 \int\mathrm DW\,e^{-\langle\bar W,OW\rangle_\nu}
     &=\det O\quad\text{for every }O.\label{eq:weighted-determinant}
\end{align}
Finally, consider a linear substitution
$\phi_i\mapsto\sum_jO_{ij}\theta_j$, and put
\[
 \|O_{ij}\|_{1,\infty}
     =\sup_t\sum_s|O_{ij}(t,s)|\nu'_j(s).
\]
If $\sum_j\|O_{ij}\|_{1,\infty}\mathsf w'_j\le\mathsf w_i$ for
each $i$, then its image $F'$ obeys
$\|F'\|_{\mathsf w'}\le\|F\|_{\mathsf w}$.
Only rows corresponding to generators occurring in $F$ need satisfy
this bound.
\end{proposition}

\begin{proof}
The identity density factors as
\[
 e^{-\langle\bar W,W\rangle_\nu}
            =\prod_t(1+\nu(t)W(t)\bar W(t)).
\]
An ordered monomial $W(J)$ contributes to its integral only if $J$
contains both generators of every mode it meets. In that case its
integral is its reordering sign times
$\prod_{(t,0)\in J}\nu(t)^{-1}$; otherwise it is zero. Its modulus
is therefore at most $\prod_{u\in J}\nu(u)^{-1/2}\le w_W^{|J|}$.
Apply this coefficientwise to obtain
\eqref{eq:weighted-integration-bound}. For disjoint mode sets a
nonzero integral requires each monomial separately to be paired.
Their degrees in $W$ are then even, so the signs from exchanging
external generators cancel. The paired-mode formula gives
factorization. It also shows that modes absent from the integrand may
be removed.

A linear change $W\mapsto OW$ multiplies the coefficient of the top
monomial by $\det O$, since the matrix of the change is
$O(t,t')\nu(t')$. Applied to the identity density this proves
\eqref{eq:weighted-determinant}, without an invertibility assumption.
For invertible $C$, use $W\mapsto CW$, leaving $\bar W$ fixed, in
the density with covariance $C$. This proves
\eqref{eq:covariance-substitution} and normalization. Its two-point
function is
\[
 \sum_s C(t,s)\nu(s)\int\dd\mu_I(W)W(s)\bar W(t')
       =\sum_s C(t,s)\nu(s)\frac{\mathbf1[s=t']}{\nu(t')}
       =C(t,t').
\]
Thus it is the kernel, rather than the matrix entry.

Under the stated substitution, the norm of the image of
$\phi_i(t)$ is at most
\[
 \sum_j\mathsf w'_j\sum_s|O_{ij}(t,s)|\nu'_j(s)\le\mathsf w_i
\]
by hypothesis. Submultiplicativity bounds
the image of each monomial by its original weight; summing its
coefficients proves the last assertion, with the same proof on the
subalgebra generated by the rows that occur.
\end{proof}

For linear combinations $v_1,\ldots,v_{2r}$ of the $W,\bar W$,
Wick's formula gives
\begin{equation}\label{eq:weighted-wick}
 \int\dd\mu_I(W)v_1\cdots v_{2r}
      =\operatorname{Pf}\left(\int\dd\mu_I(W)v_iv_j\right)_{i,j}.
\end{equation}
Odd products integrate to zero. To fix the signs, let $\partial_u$
be the left Grassmann derivative. It satisfies
$\partial_u(vG)=(\partial_uv)G-v\partial_uG$, and its Berezin
integral is zero. Differentiating the factored Gaussian density gives
\begin{align*}
 \int\dd\mu_I(W)W(t)G
      &=\nu(t)^{-1}\int\dd\mu_I(W)\partial_{\bar W(t)}G,\\
 \int\dd\mu_I(W)\bar W(t)G
      &=-\nu(t)^{-1}\int\dd\mu_I(W)\partial_{W(t)}G.
\end{align*}
Consequently, if $I(v)=\int\dd\mu_I(W)v$,
\[
 I(v_1\cdots v_q)=\sum_{j=2}^q(-1)^jI(v_1v_j)
                         I(v_2\cdots\widehat v_j\cdots v_q).
\]
This is the expansion of the Pfaffian along its first row and proves
\eqref{eq:weighted-wick} by induction. Formula
\eqref{eq:covariance-substitution} gives the corresponding Wick rule
for $\mu_C$.

In the renormalization step the modes have weight one for $k\ge1$.
At $k=0$ they are $(c,g,\beta)$ with weight $1/48$. Thus the
integration weight is
\begin{equation}\label{eq:first-step-Grassmann-weight}
 w_W=1\quad(k\ge1),\qquad w_W=\sqrt{48}\quad(k=0).
\end{equation}
In particular the first-step identity covariance has two-point
function $48\mathbf1[(c,g,\beta)=(c',g',\beta')]$.
The pair index used for a H\"older field has measure
$(\eta_k^3/48)^2$. In a substitution involving
$\delta_{\alpha,A}H_k$, its scalar component row sum is computed
with the input measure, as in
\cref{prop:weighted-Berezin-bounds}. Internal-cluster
differences are therefore controlled by the H\"older bound with the
same field weight as other pairs.

\subsection{The Gaussian block identity}\label{subsec:Gaussian-block-identity}

For this identity use a common notation for the integration space:
at $k\ge1$ let $\mathcal D_k=D_k$ and
$\mathcal A_k(A)=Q(A)$; at $k=0$ let
$\mathcal D_0=D_A+\mu$ and $\mathcal A_0(A)=Q(A)J^*$.
The integration measure is respectively counting measure or the
chamber measure $1/48$. The target is the four-component spinor
space on the lattice of spacing $L$, with measure $L^3$.
Set
\[
 \mathcal A_k^\sharp(A)=\mathcal A_k(-A)^T,\qquad
 K_k=\mathcal D_k+b_L\mathcal A_k^\sharp\mathcal A_k,
 \qquad\Gamma_k=K_k^{-1}.
\]
Thus $K_k=D_k+b_LP_L$ for $k\ge1$ and $K_0=F_0$.
Define the quadratic operator on the coarse lattice by
\[
 D_{k+1}^0=b_L-b_L^2\mathcal A_k\Gamma_k
                                       \mathcal A_k^\sharp.
\]
Let $\Phi,\bar\Phi$ be external coarse Grassmann fields and put
\[
 \Psi^c=b_L\Gamma_k\mathcal A_k^\sharp\Phi,\qquad
 \bar\Psi^c=b_L\Gamma_k^T\mathcal A_k(A)^T\bar\Phi.
\]
\begin{theorem}\label{thm:fermion-block-integration}
Assume the hypotheses of \cref{thm:chamber-locality} and
\eqref{eq:weakened-walk-contraction}, with
$\kappa_{\mathrm w}=\kappa+\kappa_0+3$ and $\kappa,\kappa_0>0$.
Let $A\in\widetilde{\mathcal G}_k$ and $0<\mu\le1$, with fermions
antiperiodic in every coordinate. Then $K_k$ is invertible.
For every Grassmann polynomial $F$ and external coarse fields
$\Phi,\bar\Phi$, the weighted integral is
\begin{multline}\label{eq:Gaussian-block-integration}
 \int\mathrm D\Psi\,
 e^{-\langle\bar\Psi,\mathcal D_k\Psi\rangle
 -b_L\langle\bar\Phi-\mathcal A_k(-A)\bar\Psi,
                  \Phi-\mathcal A_k(A)\Psi\rangle_L}
 F(\bar\Psi,\Psi)\\
 =(\det K_k)e^{-\langle\bar\Phi,D_{k+1}^0\Phi\rangle_L}
 \int\dd\mu_I(W)
 F(\bar\Psi^c+\bar W,\Psi^c+\Gamma_kW).
\end{multline}
The incoming measure is one per unit-site spin component for $k\ge1$,
and $1/48$ per chamber spin component for $k=0$; the latter
integration includes all $192$ components per photon site.
The target pairing uses measure $L^3$. The identity-Gaussian
integration weights are respectively $1$ and $\sqrt{48}$.

There are analytic functions $E_k^d(X,A)$, depending only on bonds
meeting the polymer $X$, such that
\begin{equation}\label{eq:main-determinant-quotient}
 \frac{\det K_k(A)}{\det K_k(0)}
   =\exp\sum_{X\in\mathcal P_k}
           \bigl(E_k^d(X,A)-E_k^d(X,0)\bigr),\qquad
 |E_k^d(X,A)|\le C_dM^3e^{-(\kappa+1)d_M(X)}.
\end{equation}
They are invariant under real gauge transformations, even in $A$,
and satisfy
$\overline{E_k^d(X,\overline A)}=E_k^d(X,A)$.
They are covariant under coordinate sign changes and translations
in $LM\Z^3$. The constant $C_d$ depends on the fixed
$\rho,L,\alpha,M,\kappa,\kappa_0$, and is independent of the level,
volume, connection in the stated domain and $\mu\in(0,1]$.
\end{theorem}

The block identity follows by completing the quadratic form with the
weighted transposes. Its determinant is then localized by the
resolvent integral in \cref{subsec:determinant-integral}; the polymer
estimates permit that integral to be taken term by term.

For $\Psi=\Psi^c+W'$ and $\bar\Psi=\bar\Psi^c+\bar W'$,
completion of the square is the exact identity
\begin{multline}\label{eq:block-square-completion}
 \langle\bar\Psi,\mathcal D_k\Psi\rangle
   +b_L\langle\bar\Phi-\mathcal A_k(-A)\bar\Psi,
                    \Phi-\mathcal A_k(A)\Psi\rangle_L\\
   =\langle\bar\Phi,D_{k+1}^0\Phi\rangle_L
                    +\langle\bar W',K_kW'\rangle.
\end{multline}
Indeed, write $s=b_L\mathcal A_k^\sharp\Phi$ and
$\bar s=b_L\mathcal A_k(A)^T\bar\Phi$. The left side is
\[
 \langle\bar\Psi,K_k\Psi\rangle
      -\langle\bar\Psi,s\rangle-\langle\bar s,\Psi\rangle
      +b_L\langle\bar\Phi,\Phi\rangle_L.
\]
Substitute the means $\Gamma_ks,\Gamma_k^T\bar s$. Both cross
terms cancel, and the remaining constant is
$b_L\langle\bar\Phi,\Phi\rangle_L
-\langle\bar s,\Gamma_ks\rangle
=\langle\bar\Phi,D_{k+1}^0\Phi\rangle_L$.
All factors remain in their displayed order; no conjugation of
Grassmann fields is involved.

Translation by external odd fields preserves the top coefficient
in the integration variables. Applying this to
\eqref{eq:block-square-completion} and then using
\eqref{eq:covariance-substitution} proves
\eqref{eq:Gaussian-block-integration}, with $\delta Z_k(A)=\det K_k(A)$.
At the first step this is a single integration over all chamber
components. The determinant and the integration weight are those of
\eqref{eq:level-zero-covariance} and
\eqref{eq:first-step-Grassmann-weight}; the heavy components are
included in both.

